\chapter{Zahlendarstellungen}

\section{Natürliche Zahlen}

\subsection{Ursprünge der Zahlendarstellung}
Die \textit{natürlichen Zahlen} werden so genannt, da sie auf \enquote{natürliche Weise} beim Zählen verwendet werden.
Auf dem Gebiet der heutigen Demokratischen Republik Kongo wurde im 20. Jahrhundert ein Knochen, der sogenannte \textit{Ishango-Knochen}, gefunden.
Dieser Knochen wird auf die Zeit vor etwa 20'000 Jahren vor unserer Zeitrechnung datiert.
In den Knochen wurden deutlich erkennbar von Menschen einige Kerben eingeritzt.
Die Bedeutung dieser Kerben ist unklar, es wird jedoch angenommen, dass sie eine bestimmte Anzahl festhalten sollten.

Stellen wir uns vor, dass die Kerben die Anzahl der gefangenen Fische darstellen.
Gefangene Fische durch Kerben in einem Knochen zu repräsentieren, verlangt bereits einen recht hohen Grad an
Abstraktion --- schliesslich haben gefangene Fische und Kerben in einem Knochen auf den ersten Blick keinen offensichtlichen Zusammenhang.
Weniger abstrakt wäre es, eine Darstellung zu wählen, welche direkt an das zu zählende Objekt gebunden ist.
In diesem Fall würde man also keine Kerben ritzen, sondern die entsprechende Anzahl Fische zeichnen / skizzieren.
Eine Anzahl von drei Fischen könnte also durch drei Fisch-Symbole $\twemoji{tropical fish}\twemoji{tropical fish}\twemoji{tropical fish}$
festgehalten werden.

Anstelle von Kerben in einem Knochen oder Fisch-Symbole als Markierung könnte eine bestimmte Anzahl Fische auch durch
die entsprechende Anzahl von anderen Markierungen wie Striche ($|$) oder Kieselsteine (\textbullet) repräsentiert werden:
\begin{align*}
	\text{drei Kerben} \quad\leftrightarrow\quad \twemoji{tropical fish}\twemoji{tropical fish}\twemoji{tropical fish} \quad\leftrightarrow\quad |||
	\quad\leftrightarrow\quad \text{\textbullet\textbullet\textbullet}
\end{align*}
Die Darstellung durch skizzierte Fische ist am wenigsten abstrakt, da die dabei verwendete Markierung den zu zählenden Objekten (in diesem Fall also den Fischen) direkt nachempfunden ist.
Zahlendarstellungen, welche nur ein Symbol / eine Markierung (Kerbe, \twemoji{tropical fish}, $|$, \textbullet) verwenden, werden
\tib{unäre Zahlendarstellungen}\index{Zahlendarstellung!unäre} genannt.

\begin{myremark}\phantom{1}
	\begin{enumerate}[(a)]
		\item Der Vorteil von unären Darstellungen (gegenüber \enquote{komplexeren} Darstellungen) ist ihre einfache Verständlichkeit und
		      offensichtliche Interpretation.
		\item Unäre Darstellungen sind nur für kleine natürliche Zahlen übersichtlich, da ihre Darstellungslängen linear mit der Grösse der zu
		      darstellenden Zahl wachsen (doppelt so grosse Zahl $\rightarrow$ doppelt so lange Darstellung).
		\item Unäre Darstellungen eignen sich nicht besonders gut zum Rechnen.
	\end{enumerate}
\end{myremark}

\subsection{Basisgrössen}\label{subsection:Basisgroessen}
Anstelle von nur einem einzigen Symbol, wie bei unären Darstellungen, haben bereits Kulturen der Antike unterschiedliche
\tib{Basisgrössen}\index{Basisgrössen} verwendet. \enquote{Basisgrössen} bedeuten in diesem Kontext für Symbole, die gewisse Mengen, bzw. Zahlen repräsentieren, und dies unabhängig vom Objekt, das gezählt wird. Die Römer verwendeten die sieben verschiedenen Basisgrössen:
\begin{center}
	$I$ (Eins), $V$ (Fünf), $X$ (Zehn), $L$ (Fünfzig), $C$ (Hundert), $D$ (Fünfhundert) und $M$ (Tausend).
\end{center}
\begin{myexample}\label{myexample:roman}
Die Zahl \textit{fünfhundertvierundfünfzig} könnte in den römischen Basisgrössen zum Beispiel durch
\begin{align*}
	CCCCCLIIII
\end{align*}
oder durch
\begin{align*}
	DLIIII
\end{align*}
dargestellt werden.
\end{myexample}
Wir werden uns jede Basisgrösse als einen Typ von Münze von einem bestimmten Wert vorstellen. Die darzustellende Zahl
denken wir uns dann als einen bestimmten Geldbetrag, den wir mithilfe dieser gegebenen Basisgrössen (Münztypen) auszahlen wollen. In
dieser Vorstellung hatten die Römer also ein Währungssystem von sieben verschiedene Münztypen. Der Geldbetrag \textit{einundzwanzig} könnte
dann beispielsweise durch $XXI$ (zwei Zehnermünzen und eine Einermünze) in antiker römischer Währung ausbezahlt werden. Die alten Ägypter hingegen verwendeten damals bereits Symbole (Zeichnungen) für die Basisgrössen $1, 10, 100, 1000, 10000,
	100000$ und $1000000$.

Ein Geldbetrag kann sowohl im römischen als auch ägyptischen System typischerweise auf mehrere unterschiedliche Arten
ausbezahlt werden (zum Beispiel, indem nur Einer verwendet werden). Sie kennen diese Situation vom Bezahlen mit
Bargeld --- so kann der Betrag von hundert Schweizer Franken auf viele verschiedene Arten mit Schweizer Bargeld bezahlt werden
(z.B. mit zwei Fünfzigernoten, fünf Zwanzigernoten, einer Fünfzigernote $+$ zwei Zwanzigernoten $+$ zwei Fünfliber, etc.). Natürlich ist
es bei Zahlendarstellungen vorteilhaft, wenn sie \textit{eindeutig} sind. Dies bedeutet, dass es für jede natürliche Zahl eine einzige minimale Münzenmenge gibt, die diese Zahl repräsentiert. Angenommen wir haben einen Betrag $x$ gegeben, welcher durch
unterschiedliche Ansammlungen von Münzen ausbezahlt werden kann. Dann betrachten wir diejenige Ansammlung von Münzen als
die \tib{Zahlendarstellung}\index{Zahlendarstellung} von $x$, welche am wenigsten Münzstücke benötigt.

\begin{myexample}\label{myexample:notUnique}
Der Betrag \textit{sieben} kann im römischen System auf genau zwei verschiedene Arten (abgesehen von der Reihenfolge) ausbezahlt werden:
\begin{itemize}
	\item $IIIIIII$ (benötigt sieben Münzstücke),
	\item $VII$ (benötigt drei Münzstücke).
\end{itemize}
Die Darstellung $VII$ verwendet dabei am wenigsten Münzstücke und wird somit als die \textit{minimale} Zahlendarstellung von
\textit{sieben} verstanden.
\end{myexample}
Nun gibt es aber denkbare Münzsysteme, in denen nicht jeder Betrag auf eindeutige Weise mit möglichst wenigen
Münzstücken ausbezahlt werden kann. Dies werden Sie in \cref{myexercise:1001} untersuchen.

\begin{myexercise}\label{myexercise:1001}
	Erweitern Sie das römische Münzsystem um zwei weitere Münztypen Ihrer Wahl und finden Sie einen Betrag, der in Ihrem
	neuen System mehrere unterschiedliche Darstellungen mit derselben minimalen Anzahl von Münzen erlaubt.
	\begin{myanswer}
		Wir erweitern das römische Münzsystem um die beiden Münztypen \coin{II} und \coin{III}. Dann besitzt der Betrag
		\textit{vier} die beiden Darstellungen \coin{I}\coin{III} oder \coin{II}\coin{II}, welche beide eine (minimale) Anzahl von zwei Münzstücken verwenden.
	\end{myanswer}
\end{myexercise}


\begin{myexercise}\label{myexercise:grosse-roem-zahlen}
	Ein Haupt-Nachteil von Zahlendarstellungen durch Münzsysteme ist, dass Münzsysteme für grosse Zahlen zur gleichen Ineffizienz führen wie unäre Zahlensysteme. Wie viele Zeichen würde man benötigen, um die Zahl 1'000'000'000 (1 Milliarde) im römischen Zahlensystem darzusetellen?
	\begin{myanswer}
		Die Zahl wäre eine Million Stellen lang: Die Zahl \textit{M} (Tausend) ist die grösste Zahl des römischen Zahlensystems und wir benötigen eine Million mal Tausend, um auf eine Milliarde zu kommen.
	\end{myanswer}
\end{myexercise}
\subsection{Stellenwertdarstellung}
\label{subsec:stellenwertdarst}
Das uns vermutlich am besten vertraute Zahlensystem ist das \tib{Dezimalsystem}\index{Stellenwertsystem!Dezimalsystem}
(Zehnersystem). Es verwendet die \textbf{zehn} verschiedenen Ziffern
\begin{align*}
	\lrc{0,1,2,3,4,5,6,7,8,9}
\end{align*}
und die (natürlichen) Potenzen $10^0, 10^1, 10^2, \ldots$ der Basis $10$ als \tib{Stellenwerte}\index{Stellenwert}.
Bereits in der Primarschule haben Sie gelernt, dass beispielsweise das Symbol $\green{4}\orange{6}\blue{0}\red{5}$ definiert ist als die folgende
Summe:
\begin{center}
	\red{fünf} Einer ($\red{5}\cdot 10^0$) $+$ \blue{null} Zehner ($\blue{0}\cdot 10^1$) $+$ \orange{sechs} Hunderter
	($\orange{6}\cdot 10^2$) + \green{vier} Tausender
	($\green{4}\cdot 10^3$).
\end{center}
Eine Ziffer gibt dabei jeweils an, wie viele Male die entsprechende Potenz der Basis $10$ verwendet werden soll. Das
Dezimalsystem hat also (die unendlich vielen verschiedenen) Basisgrössen
\begin{align*}
	10^0=1,\quad 10^1=10,\quad 10^2=100,\quad 10^3=1000,\ldots
\end{align*}

\begin{mydefinition}\label{mydefinition:Stellenwertdarstellung}
Seien $a_0, a_1, \ldots, a_n$ natürliche Zahlen. Dann ist der Ausdruck
\begin{align*}
	a_na_{n-1}\ldots a_1a_0
\end{align*}
definiert durch die Summe
\begin{align*}
	a_n\cdot 10^n + a_{n-1}\cdot 10^{n-1}+\ldots + a_1\cdot 10^1 + a_0\cdot 10^0.
\end{align*}
\end{mydefinition}

\begin{myexample}
Der Ausdruck $4205$ entspricht der Zahl
\begin{align*}
	\textcolor{red}{4205} := \textcolor{red}{\overbrace{4}^{a_3}}\cdot \underbrace{10^3}_{=1000} + \textcolor{red}{\overbrace{2}^{a_2}}\cdot \underbrace{10^2}_{=100} + \textcolor{red}{\overbrace{0}^{a_1}}\cdot \underbrace{10^1}_{=10} + \textcolor{red}{\overbrace{5}^{a_0}}\cdot \underbrace{10^0}_{=1}.
\end{align*}
\end{myexample}

Anstelle der Potenzen der Basis $10$ kann man auch andere natürliche Zahlen als Basis verwenden.
Das \tib{Binärsystem}\index{Stellenwertsystem!Binärsystem} (Zweiersystem) verwendet zum Beispiel die Basis $2$ und
Ziffern aus $\lrc{0,1}$.
\begin{myexample}
Wir wollen den Ausdruck $1101$ in Basis $\blue{2}$ interpretieren. Um dies zu verdeutlichen, schreibt man die Basis
tiefgestellt dahinter, also $1101_{\blue{2}}$. Der Ausdruck $1101_2$ wird dann aufgefasst als die Zahl
\begin{align*}
	1\cdot 2^3 + 1\cdot 2^2 + 0\cdot 2^1 + 1\cdot 2^0 = 1\cdot 8 + 1\cdot 4 + 0\cdot 2 + 1\cdot 1 = 13.
\end{align*}
Die binäre Zahl $1101_2$ entspricht also der Zahl dreizehn.
\end{myexample}

\begin{myattention}
	Falls wir nicht explizit eine Basis tiefgestellt hinter einen Zahlenausdruck schreiben, dann ist dieser Ausdruck immer in Basis $10$ (im
	Dezimalsystem) aufzufassen. Somit meinen wir also mit $11$ die Zahl elf und nicht etwa die Zahl $11_2 = 1\cdot 2^1 +
		1\cdot 2^0$.
\end{myattention}

\begin{myexercise}
	Schreiben Sie die Binärzahl $11011_2$ als Dezimalzahl.
	\begin{myanswer}
		\begin{align*}
			11011_2 = 1\cdot 2^4 + 1\cdot 2^3 + 0\cdot 2^2 + 1\cdot 2^1 + 1\cdot 2^0 = 16 + 8 + 2 + 1 = 27
		\end{align*}
	\end{myanswer}
\end{myexercise}

\begin{myexercise}
	Schreiben Sie die Dezimalzahl $14$ als Binärzahl.
	\begin{myanswer}
		\begin{align*}
			14 = 8 + 4 + 2 = 1\cdot 2^3 + 1\cdot 2^2 + 1\cdot 2^1 + 0\cdot 2^0 = 1110_2
		\end{align*}
	\end{myanswer}
\end{myexercise}

Warum schränken wir uns beim Schreiben von Binärzahlen eigentlich auf den Gebrauch von lediglich den beiden Ziffern aus
$\lrc{0,1}$ ein? Ein Ausdruck der Form
\begin{align*}
	3\cdot 2^2 + 5\cdot 2^1 + 7\cdot 2^0
\end{align*}
wäre doch ebenfalls sinnvoll. Diese Frage wollen wir in den folgenden Aufgaben klären. Erinnern Sie sich zurück an \cref{subsection:Basisgroessen}.
Dort haben wir versucht, Geldbeträge mit möglichst wenigen Münzstücken auszuzahlen.

\begin{myexercise}
	Der Ausdruck $3\cdot 2^2 + 5\cdot 2^1 + 7\cdot 2^0$ verwendet $3+5+7 = 15$ Münzstücke und stellt die Zahl $29$ dar.
	Stellen Sie die Zahl $29$ im Binärsystem dar, verwenden Sie diesmal aber lediglich Ziffern aus $\lrc{0,1}$. Wie viele
	Münzstücke verwendet diese Darstellung?
	\begin{myanswer}
		\begin{align*}
			29 = 1\cdot 2^4 + 1\cdot 2^3 + 1\cdot 2^2 + 0\cdot 2^1 + 1\cdot 2^0
		\end{align*}
		Diese Darstellung verwendet lediglich 5 Münzstücke.
	\end{myanswer}
\end{myexercise}

\begin{myexample}\label{myexample:ersetzen}
Wir wollen dreizehn Objekte im Dezimalsystem mit so wenigen Münzstücken wie möglich darstellen. Wir könnten dies mit
dreizehn Einern tun. Doch wir sehen sofort, dass sich zehn Einer stets durch nur einen einzigen Zehner (und null
Einer) ersetzen lassen. Dadurch haben wir insgesamt neun Münzen gespart:
\begin{itemize}
	\item verfügbare Münztypen: \coin{1}, \coin{10}, $\ldots$
	\item Münzen vorher:\\
	      \coin{1}\coin{1}\coin{1}\coin{1}\coin{1}\\
	      \coin{1}\coin{1}\coin{1}\coin{1}\coin{1}\\
	      \coin{1}\coin{1}\coin{1}\\
	\item Münzen nachher:\\
	      \coin{10}\coin{1}\coin{1}\coin{1}\\
\end{itemize}
\end{myexample}

\begin{mydefinition}[$\blue{b}$-adische Zahlendarsellung für natürliche Zahlen]
\label{mydefinition:badisch}
Gegeben ist eine natürliche Zahl $\blue{b}\geq 2$ (\tib{Basis}\index{Basis}). Dann definieren wir
\begin{align*}
	(\red{a}_{n}\red{a}_{n-1}\ldots \red{a}_{1}\red{a}_{0})_{\blue{b}} := \red{a}_n\cdot \blue{b}^n + \red{a}_{n-1}\cdot \blue{b}^{n-1} + \ldots +\red{a}_1\cdot \blue{b}^1 + a_0\cdot
	\blue{b}^0,
\end{align*}
wobei die Zahlen $a_{0}, a_{1},\ldots, a_{n-1}, a_{n}\in\lrc{0,1,\ldots,b-1}$ \tib{Ziffern}\index{Ziffern} genannt werden. Wird
nicht explizit eine andere Basis genannt, so ist stets die Basis $b=10$ gemeint.

Diese Zahlendarstellung wird $b$-\tib{adische Zahlendarstellung} genannt.
\end{mydefinition}

\begin{myexample}
Für die Wahl $b:=3$ erhalten wir die $3$-adische Zahlendarstellung (Dreiersystem). Die Zahl $211_3$ entspricht dann der Zahl
\begin{align*}
	211_3 = 2\cdot 3^2 + 1\cdot 3^1 + 1\cdot 3^0 = 2\cdot 9 + 1\cdot 3 + 1\cdot 1 = 22.
\end{align*}
\end{myexample}

\begin{mychallenge}
	Begründen Sie, warum es nicht besonders sinnvoll ist, die natürliche Zahl $1$ als Basis zu wählen, selbst wenn wir in \cref{mydefinition:badisch} die Einschränkung $a_0, a_1, \ldots, a_{n-1}, a_{n}\in\lrc{0,1,\ldots,b-1}$ an die Ziffern ignorieren?
	\begin{myanswer}
		Mit der Wahl der Basis $b := 1$ erhalten wir keine unterschiedlichen Stellenwerte, da $b^k = 1^k = 1$ für alle natürlichen Zahlen $k$ gilt.
	\end{myanswer}
\end{mychallenge}

\begin{myremark}
	Warum ist es sinnvoll in \cref{mydefinition:badisch} zu fordern, dass eine $b$-adische Darstellung keine Ziffern grösser als $b-1$ verwenden darf?
	Unsere Argumentation ist ähnlich zu der in \cref{myexample:ersetzen}. Angenommen eine Darstellung
	\begin{align*}
		\red{a}_n\cdot \blue{b}^n + \red{a}_{n-1}\cdot \blue{b}^{n-1} + \ldots +\red{a}_1\cdot \blue{b}^1 + a_0\cdot
		\blue{b}^0
	\end{align*}
	würde eine Ziffer $a_k \geq b$ verwenden. Anders gesagt, würden dann $b$ oder mehr Münzstücke von dem Typ $b^k$ verwendet werden.
	Dann könnten wir aber immer eine Gruppe von $b$ vielen Münzen vom Typ $b^k$ durch jeweils eine einzige Münze vom Typ $b^{k+1}$ ersetzen und würden für jede solche Gruppe $b-1$ Münzstücke einsparen.

	Beispielsweise ist es in der Basis $b = 10$ nicht sinnvoll 23 Hunderter ($23\cdot 10^2$) zu verwenden, da wir stattdessen zwei Tausender und nur drei Hunderter ($2\cdot 10^3 + 3\cdot 10^2$) verwenden könnten.
	Dabei würden wir insgesamt $18$ Münzstücke einsparen.
\end{myremark}

\begin{myexercise}
	Ordnen Sie die folgenden Zahlen aufsteigend:
	\begin{align*}
		43_{10}, 2201_3, 11101_2, 33_4, 142_5, 111_2.
	\end{align*}
	\begin{myanswer}
		\begin{align*}
			2201_3  & = 2\cdot 3^3 + 2\cdot 3^2 + 0\cdot 3^1 + 1\cdot 3^0 = 73,              \\
			11101_2 & = 1\cdot 2^4 + 1\cdot 2^3 + 1\cdot 2^2 + 0\cdot 2^1 + 1\cdot 2^0 = 29, \\
			33_4    & = 3\cdot 4^1 + 3\cdot 4^0 = 15,                                        \\
			142_5   & = 1\cdot 5^2 + 4\cdot 5^1 + 2\cdot 5^0 = 47,                           \\
			111_2   & = 1\cdot 2^2 + 1\cdot 2^1 + 1\cdot 2^0 = 7,                            \\
			        & \Rightarrow 111_2 < 33_3 < 11101_2 < 43 < 142_5 < 2201_3
		\end{align*}
	\end{myanswer}
\end{myexercise}

\begin{mychallenge}
	Die nachfolgende Aussage gilt für jede Basis $b$ eines $b$-adischen Systems. Zur Vereinfachung werden wir aber nur die Wahl der Basis $b := 10$ untersuchen.
	Zeigen Sie, dass die $10$-adische Darstellung eindeutig ist. Es gibt also keine zwei unterschiedlichen $10$-adischen Darstellungen $x$ und $y$, welche dieselbe Zahl (Anzahl Objekte) darstellen.
	\begin{myanswer}
		Seien
		\begin{align*}
			x & := a_n\cdot 10^n + a_{n-1}\cdot 10^{n-1} + \ldots + a_k\cdot 10^k + a_{k-1}\cdot 10^{k-1} + \ldots + a_1\cdot 10^1 + a_0\cdot 10^0 \\
			y & := b_n\cdot 10^n + b_{n-1}\cdot 10^{n-1} + \ldots + b_k\cdot 10^k + b_{k-1}\cdot 10^{k-1} + \ldots + b_1\cdot 10^1 + b_0\cdot 10^0
		\end{align*}
		zwei unterschiedliche $10$-adische Darstellungen derselben Zahl.
		Dann müssen sich $x$ und $y$ an mindestens einer Stelle (in einer Ziffer) unterscheiden.
		Sei $10^k$ die grösste Stelle, an der sich $x$ und $y$ unterscheiden. Wir können annehmen, dass $a_k > b_k$ gilt. Dann bewirkt der Unterschied an dieser Stelle
		eine Differenz von mindestens $10^k$. Doch selbst mit der Wahl $a_{k-1} = \ldots = a_1 = a_0 = 0$ und $b_{k-1} = \ldots = b_1 = b_0 = 9$ kann nur noch ein Unterschied von $10^k - 1$ erwirkt werden.
		Dann ist aber die von $x$ dargestellte Zahl um mindestens 1 grösser als die von $y$ dargestellte Zahl.
	\end{myanswer}
\end{mychallenge}

\subsection{Umrechnungen zwischen Darstellungen}
Wir wollen nun sehen, wie verschiedene Darstellungen ineinander umgerechnet werden können.
\begin{myexample}[Darstellung in Basis $5$]\label{myexample:umrechnung}
Wie lautet die $5$-adische Darstellung der Zahl $84$? Wir beginnen damit alle natürlichen Potenzen von $5$
aufzulisten, die kleiner oder gleich $84$ sind:
\begin{align*}
	5^0 = 1, \quad 5^1 = 5, \quad 5^2 = 25.
\end{align*}
Die Potenz $5^3 = 125 > 84$ ist bereits grösser als die gegebene Zahl. Nun versuchen wir möglichst viele der möglichst
grössten Potenz von $5$ zu verwenden, solange diese noch in $84$ passen. Die grösste mögliche Potenz ist $5^2 = 25$ und
$25$ hat \red{dreimal} in $84$ Platz. Es bleibt also noch $84-3\cdot 5^2 = 9$ übrig. Wir verwenden möglichst viele der
möglichst grössten Potenz von $5$, die in $9$ Platz haben und stellen fest, dass $5^1=5$ \blue{einmal} Platz hat. Dann
bleibt noch $9-5^1 = 4$ übrig. Wir verwenden möglichst viele der möglichst grössten Potenz von $5$, die in $4$ Platz
haben. Wir verwenden also \green{viermal} die Potenz $5^0 = 1$. Nun bleibt nichts mehr übrig und die $5$-adische
Darstellung von $84_{10}$ lautet:
\begin{align*}
	84_{10} = \red{3}\blue{1}\green{4}_5.
\end{align*}
\end{myexample}
Da wir in \cref{myexample:umrechnung} immer mit der grössten Basisgrösse begonnen haben und so viele Münzen wie möglich
davon verwendet haben, nennen wir diese Methode \tib{gierig}\index{gierig / greedy} (englisch: \tib{greedy}).

\subsection{Hexadezimalsystem}
Wie Sie in \cref{subsec:stellenwertdarst} gesehen haben, kann jede Zahl im Grunde genommen mit (fast) beliebig vielen Ziffern dargestellt werden. Während das Binärsystem die Ziffern $\lrc{0,1}$ verwendet, verwenden wir im (uns geläufigeren) Dezimalsystem die Ziffern $\lrc{0,1,2,3,4,5,6,7,8,9}$. Wir haben zudem auch beobachtet, dass die höchste Ziffer in einem Zahlensystem mit $b$ Ziffern immer $b-1$ ist: Beispielsweise ist im Binärsystem, das 2 Ziffern hat, die höchste Ziffer eine 1 ($= 2-1$). Im Dezimalsystem, das bekanntlich 10 Ziffern hat, ist die höchste Ziffer eine 9 ($ = 10-1$). Die höchste Ziffer in einem $b$-adischen System ist also immer $b-1$.

Welche Ziffern sollen wir jedoch wählen, sofern $b>10$ gewählt wird? Die dezimale Zahl 10 ist ja keine Ziffer, sondern besteht selber aus zwei Ziffern (genau so, wie es das Dezimalsystem vorschreibt).

Im Grunde genommen sind Ziffern wie \enquote{1}, \enquote{2} usw. nichts weiter als \enquote{Zeichnungen}, die von Menschen angefertigt worden sind, und mit welchen eine gewisse Menge, bzw. Anzahl gemeint ist. Wir könnten also einfach andere, uns bekannte \enquote{Zeichnungen} verwenden, um höhere Mengen zu bezeichnen, also etwa Buchstaben. Genau dies wird im sogenannten \tib{Hexadezimalsystem}\index{Hexadezimalsystem} gemacht ($b=16$). Die möglichen Zeichen im Hexadezimalsystem sind:
\begin{align*}
	\lrc{0,1,2,3,4,5,6,7,8,9,A,B,C,D,E,F},
\end{align*}
wobei beispielsweise die Ziffer $B$ die (dezimale) Zahl $11$ (Elf) bezeichnet.

Zur Umrechnung von Zahlen vom einen in das andere System verfolgen wir denselben Ansatz wie in \cref{myexample:umrechnung}

\begin{myexercise}
	Wie lautet die Zahl $27_{10}$ in der $16$-adischen Darstellung?
	\begin{myanswer}
		\begin{align*}
			27_{10} = 1\cdot 16^1 + 11\cdot 16^0 = 1B_{16}
		\end{align*}
	\end{myanswer}
\end{myexercise}

\begin{myexercise}
	Geben Sie folgende Zahlen im hexadezimalen System aus: 44, 23, 90, 124, 306.
	\begin{myanswer}
		2C, 17, 5A, 7C, 132
	\end{myanswer}
\end{myexercise}


\subsection{Rechnen im Binärsystem}

Rechnen im Binärsystem funktioniert genauso, wie Sie es für das Dezimalsystem gelernt haben. Beispielsweise, wenn Sie die Zahlen $49035_{10}$ und $718293_{10}$ addieren:

\[
	\begin{array}{ccccccc}
		  & 0              & 4               & 9              & 0               & 3              & 5_{10}              \\
		+ & \underset{}{7} & \underset{1}{1} & \underset{}{8} & \underset{1}{2} & \underset{}{9} & \underset{}{3_{10}} \\
		\hline
		  & 7              & 6               & 7              & 3               & 2              & 8_{10}              \\\hline\hline
	\end{array}
\]

Die kleinen Ziffern in der Rechnung oben stellen den Übertrag\index{Übertrag} (Englisch: \textit{carry over} oder einfach nur \textit{carry}) dar. Ähnlich dazu gehen Sie im Binärsystem vor. Dabei müssen im Wesentlichen nur folgende 5 Regeln beachtet werden:

\begin{itemize}
	\item 0 + 0 = 0
	\item 0 + 1 = 1
	\item 1 + 0 = 1
	\item 1 + 1 = 0, Übertrag 1
	\item 1 + 1 + 1 = 1, Übertrag 1
\end{itemize}

Ein Beispiel sei unten angegeben:
\[
	\begin{array}{ccccccc}
		  & 1               & 0               & 1               & 0               & 1               & 1_2              \\
		+ & \underset{1}{1} & \underset{1}{1} & \underset{0}{1} & \underset{1}{0} & \underset{0}{1} & \underset{}{0_2} \\
		\hline
		1 & 1               & 0               & 0               & 1               & 0               & 1_2              \\\hline\hline
	\end{array}
\]

\begin{myexercise}
	Rechnen Sie folgende beiden Zahlen im Binärsystem zusammen: $10001010_2$, $11101_2$.
	\begin{myanswer}
		$10100111_2$
	\end{myanswer}
\end{myexercise}

\begin{myexercise}
	Welches ist die grösste 4-stellige Binärzahl (in Dezimalschreibweise)?
	\begin{myanswer}
		\begin{align*}
			1111_2 = 2^4 + 2^3 + 2^2 + 2^1 = 8 + 4 + 2 + 1 = 15_{10}
		\end{align*}
	\end{myanswer}
\end{myexercise}

\begin{mychallenge}
	Es sei $n>0$ eine natürliche Zahl. Es sei $z$ die grösste $n$-stellige Binärzahl. Wie sieht $z$ aus? Finden Sie eine möglichst einfache Darstellung von $z$.
	\begin{myanswer}
		Die grösste $n$-stellige Binärzahl $z$ ist gegeben durch
		\begin{align*}
			z := \underbrace{1\ldots 11_2}_{n\text{-Stellen}} = 1\cdot 2^0 + 1\cdot 2^1 + 1\cdot 2^2 + \ldots + 1\cdot 2^{n-1}.
		\end{align*}
		Die Umrechnung in Basis 10 ist aufwendig, da wir die Summe $2^0 + 2^1 + 2^2 + \ldots + 2^{n-1}$ berechnen müssten. Lassen Sie uns die Zahl $z+\textcolor{Red}{1}$ bestimmen:
		\[
			\begin{array}{ccccccc}
				  & 1               & 1               & \ldots & 1               & 1               & 1_2                               \\
				+ & \underset{1}{0} & \underset{1}{0} & \ldots & \underset{1}{0} & \underset{1}{0} & \underset{}{\textcolor{Red}{1}_2} \\
				\hline
				1 & 0               & 0               & \ldots & 0               & 0               & 0_2                               \\\hline\hline
			\end{array}
		\]
		Wir haben die kleinste $n+1$-stellige Binärzahl erhalten. Wir sehen nun, dass $z+1$ geben ist durch
		\begin{align*}
			z+1 = \underbrace{10\ldots 00_2}_{n+1\text{-Stellen}} = 0\cdot 2^0 + 0\cdot 2^1 + \ldots + 0\cdot 2^{n-1} + 1\cdot 2^n = 2^n.
		\end{align*}
		Es fallen also alle Summanden bis auf den Term $ 1\cdot 2^n$ weg. Da $z$ genau $1$ kleiner ist als $z+1$, finden wir $z = 2^n - 1$.
	\end{myanswer}
\end{mychallenge}

\begin{myexercise}
	Rechnen Sie folgende beiden Zahlen im Binärsystem zusammen: $11110001_2$, $00001111_2$.
	\begin{myanswer}
		Die mathematisch richtige Antwort wäre eigentlich $100000000_2$. Falls wir die Rechnung jedoch in einer 8-Bit-Umgebung ausführen (wie es die ersten Computer waren), könnte das Resultat auch 0 sein, da die Zahl zu lang würde (Überlauf, en. \textit{overflow}).
	\end{myanswer}
\end{myexercise}

\begin{mychallenge}\label{challenge:online}
	\begin{itemize}
		\item Wie lange benötigen Sie für 6 Umwandlungen? Testen Sie sich \href{https://oinf.ch/interactive/umrechnen-zwischen-zahlensystemen/}{hier}.
		\item Testen Sie Ihre Zahlensystem-Umwandlungs-Wissen \href{https://oinf.ch/interactive/basiswechsel-10-zu-2/}{hier}!
		\item Schauen Sie sich das Video der Die \href{https://www.youtube.com/watch?app=desktop&v=GcDshWmhF4A&embeds_referring_euri=https%3A%2F%2Fwww.yaplakal.com%2Fforum28%2Fst%2F125%2Ftopic930979.html&feature=emb_imp_woyt}{Marble Adding Machine} an.
		\item Schauen Sie sich an, wie man einen binären Computer auch \href{https://youtu.be/5wKz1e0ULcg?si=AiPXUniyfU15X11m}{mit Wasser} bauen kann.
	\end{itemize}
\end{mychallenge}
